\section{Relations and functions}
\label{sec:set-prerequisites-03}

\begin{definition}[Binary relation]
A \emph{binary relation from \(A\) to \(B\)} is a subset of \(A\times B\).  A
\emph{binary relation on \(A\)} is a subset of \(A\times A\).
\end{definition}

\begin{definition}[Properties of a binary relation]
Let \(R\) be a binary relation on \(A\).  It is \emph{reflexive} iff
\(aRa\) for every \(a\in A\); \emph{irreflexive} iff \(aRa\) for no
\(a\in A\); \emph{symmetric} iff \(aRb\) implies \(bRa\) for all
\(a,b\in A\); and \emph{transitive} iff \(aRb\) and \(bRc\) imply \(aRc\)
for all \(a,b,c\in A\).

The relation \(R\) is a \emph{strict linear order} on \(A\) iff it is
irreflexive and transitive and, for every distinct \(a,b\in A\), exactly one
of \(aRb\) and \(bRa\) holds.  It \emph{strictly well-orders} \(A\) iff it is
a strict linear order and every nonempty \(B\subseteq A\) has an
\(R\)-least member: some \(b\in B\) such that no \(a\in B\) satisfies
\(aRb\).  These clauses define the relation properties used below; none is a
second name for another clause, and strict well-ordering is required in the
definition of ordinal.
\end{definition}

\begin{definition}[Function]
A \emph{function from \(A\) to \(B\)}, written \(f:A\to B\), is a binary
relation \(f\subseteq A\times B\) such that for every \(a\in A\) there exists
exactly one \(b\in B\) with \(\tuple{a,b}\in f\).  The set \(A\) is the
domain, and \(B\) is the codomain.
\end{definition}

The value at \(a\) is written \(f(a)\).  The range is
\(\rng(f)=\set{f(a):a\in A}\).  The codomain is part of the typed declaration
and need not equal the range.

\begin{definition}[Partial function]
A \emph{partial function from \(A\) to \(B\)}, written
\(f:A\rightharpoonup B\), is a function \(f:D\to B\) for some
\(D\subseteq A\).  Its domain of definition is \(D\).
\end{definition}

\begin{definition}[Injection, surjection, and bijection]
Let \(f:A\to B\).
\begin{enumerate}
  \item The function is \emph{injective} iff
        \(f(a)=f(a')\) implies \(a=a'\) for all \(a,a'\in A\).
  \item The function is \emph{surjective} iff for every \(b\in B\) there is
        an \(a\in A\) such that \(f(a)=b\).
  \item The function is \emph{bijective} iff it is injective and surjective.
\end{enumerate}
\end{definition}

\begin{definition}[Image and inverse image]
For \(X\subseteq A\) and \(Y\subseteq B\), define
\[
  f[X]\defeq\set{f(x):x\in X},
  \qquad
  f^{-1}[Y]\defeq\set{x\in A:f(x)\in Y}.
\]
The notation \(f^{-1}[Y]\) does not assert that \(f\) has an inverse function.
\end{definition}

